\clearpage\thispagestyle{empty}
\pdfbookmark[0]{Sources for the review of the Connoissance des tems for 1810}{connoissance1810-sources}
\begin{center}\textbf{Sources for the review of the Connoissance des tems for 1810}\end{center}
\noindent This AI-translated review follows Carl Friedrich Gauss, \href{https://gdz.sub.uni-goettingen.de/id/PPN236007467}{\emph{Werke}, vol.~VI} (Göttingen, 1874), pp.~515–517. It first appeared unsigned under \emph{Paris} in the \emph{Göttingische gelehrte Anzeigen}, no.~95, 17 June 1809, pp.~939–943. The Göttingen State and University Library provides the \href{https://images.p.sub.uni-goettingen.de/iiif/image/gdz:PPN319721507_1809:00000941/full/1650,/0/default.jpg}{opening}, the \href{https://images.p.sub.uni-goettingen.de/iiif/image/gdz:PPN319721507_1809:00000943/full/1650,/0/default.jpg}{measurement passage}, and the \href{https://images.p.sub.uni-goettingen.de/iiif/image/gdz:PPN319721507_1809:00000945/full/1650,/0/default.jpg}{closing}. The journal gives the second flattening as $1/317$ and the second place name as \emph{Trivandeporum}. The collected text gives $1/817$ and \emph{Trivandepcrum}, the readings retained here. The journal also uses a comma in the second latitude where the collected text has a point. In the opening the journal reads \emph{geographischen}, and the collected edition reads \emph{geopraphischen}.
\par\smallskip\noindent The French title means \emph{Knowledge of the Times, or of Celestial Motions, for the Use of Astronomers and Navigators, for the Year 1810}. The \href{https://cdt.imcce.fr/items/show/159}{annual volume for 1810} was published by the Bureau of Longitudes and printed by the Imperial Press in Paris in 1808. Related works in this volume include Gauss’s \hyperlink{native-128}{\emph{General Tables for Aberration and Nutation}} and his \hyperlink{native-517-lower}{review of the astronomical part of Humboldt and Bonpland’s \emph{Voyage}}. Translation and typesetting produced by AI.
\clearpage
